\subsection{Functors}

DEFINITION 3. --- Let C and C' be categories. A functor from C to C' is called a morphism of quivers \(\varphi\colon\mathrm{C}\to\mathrm{C}'\) such that \(\varphi(fg)=\varphi(f)\varphi(g)\) for every pair \((f,g)\) of composable arrows of C.

Let C and C' be categories and let \(\varphi\colon\mathrm{C}\to\mathrm{C}'\) be a functor. Let \(f\) be an arrow of C, \(a\) its origin and \(b\) its target; then, \(\varphi(f)\) is an arrow of C' with origin \(\varphi(a)\) and target \(\varphi(b)\) .

Let C, C', C'' be categories and let \(\varphi\colon\mathrm{C}\to\mathrm{C}'\) and \(\varphi'\colon\mathrm{C}'\to\mathrm{C}''\) be functors. Then, \(\varphi'\circ\varphi\) is a functor.

Let C be a category. Then the morphism of quivers \(Id_{C}\) is a functor.

Let \(\varphi\colon\mathrm{C}\to\mathrm{C}^{\prime}\) be a functor. In order for \(\varphi\) to be an isomorphism, it is necessary and sufficient that the maps \(\operatorname{Som}(\varphi)\) and \(\operatorname{Fl}(\varphi)\) be bijective.

Thus, if one calls a category structure on sets S and F the data of a quiver structure on these sets and of an associative composition law for which there exists at each vertex an identity element, one can take functors as morphisms of the category structure (E, IV, p. 11).

